% !TeX root = ../main.tex

\newgeometry{left=1cm, right=1cm, top=1cm, bottom=1cm}

\name{Nils Jörgensen}
\tagline{Autonomous Systems Researcher | Motion \& Mission Planning | Industrial Robotics}
\photoR{2.5cm}{picture}
\personalinfo{%
  \mailaddress{Sofielundsvägen 5, 121 32 Enskededalen, Sweden}
  \phone{070-581 97 06}
  \email{nilsjor@kth.se}
}

\makecvheader

%% Set the left/right column width ratio
\columnratio{0.615}

\begin{paracol}{2}

%% ---- LEFT COLUMN -------------------------------------------------------

\cvsection{Experience}

\cvevent{Ph.\,D. Researcher | Networked Industrial Robotics}{KTH Royal Institute of Technology}{Autumn 2020 -- Spring 2026}{Stockholm, Sweden}
Research spanning task-, motion-, and mission planning for safe industrial multi-robot coordination, including PDDL, AI-planning, and MILP optimization. Surveyed planning algorithms for indoor robotics and developed joint planners integrating task scheduling with wireless network resource allocation, focusing on logistics.

\divider

\cvevent{Teacher, Student Supervisor, and Research Engineer}{KTH | Department of Machine Design}{Autumn 2019 -- Spring 2026}{Stockholm, Sweden}
Years of teaching, course development, and research, covering computer architectures, assembly and C/C\texttt{++} bare-metal and RTOS programming, distributed systems, modeling and control of dynamical systems, and everything in between.
Supervised countless in degree projects as well as industrial proof-of-concept projects within the Mechatronics capstone course.

\hfill

\cvsection{Education}

\cvevent{Ph.\,D. in Mechatronics and Embedded Control Systems}{KTH Royal Institute of Technology}{Autumn 2020 -- Spring 2026}{Stockholm, Sweden}

Formal Methods \& Software Verification \textbullet{}
Hybrid Systems \textbullet{}
Networked Control Systems \textbullet{}
Cyber-Physical Security \textbullet{}
Reinforcement Learning \textbullet{}
Convex Optimization \textbullet{}
Fault-Tolerant Design \textbullet{}
Research Methodology \textbullet{}
Scientific Writing

\divider

\cvevent{M. Sc. in Engineering Design | Mechatronics Track}{KTH Royal Institute of Technology}{Autumn 2017 -- Spring 2019}{Stockholm, Sweden}

Internetworking \textbullet{} 
Internet Security and Privacy \textbullet{} 
Embedded and Distributed Control Systems \textbullet{} 
Nonlinear Control \textbullet{}
Optimal and Model-Predictive Control \textbullet{}
Low-Level Programming \textbullet{}
Robust Electronics \textbullet{}
Test-Driven Development \textbullet{} 
Systems Engineering

\divider

\cvevent{B. Sc. in Mechanical Engineering | International Program}{RWTH Aachen University / KTH Royal Institute of Technology}{Autumn 2014 -- Summer 2017}{Aachen, Germany}

Automatic Control \textbullet{}
Object-Oriented Programming \textbullet{}
Algorithms and Optimisation \textbullet{}
Statistics and Probability Theory \textbullet{}
Numerical Analysis \textbullet{}
Machine Components \textbullet{}
Electrical Engineering \textbullet{}
Physics

%% ---- RIGHT COLUMN ------------------------------------------------------
\switchcolumn

\cvsection{Languages}

\cvskill{Swedish}{5}
\cvskill{English}{5}
\cvskill{German}{4}

\hfill

\cvsection{Technical Skills}

\cvachievement{\faCogs}{Planning \& optimization}{Mission \& motion planning with a focus on PDDL, MILP, and other optimization-based methods.}

\divider

\cvachievement{\faSitemap}{Networked multi-robot systems}{Multi-robot task allocation, and QoS-aware replanning in dynamic environments. Joint optimization of task and resource scheduling.}

\divider

\cvachievement{\faCheckSquare[regular]}{Formal methods \& systems theory}{Model checking, satisfiability solving, hybrid-systems analysis. Nonlinear and optimal control.}

\divider

\cvachievement{\faCode}{C++, Python, MATLAB and ROS2}{Modeling, simulation, prototyping and deployment: bare-metal, RTOS, and application-level development.}

\hfill

\cvsection{Publications}

\cvachievement{\faFile*[regular]}{Joint 5G-aware task planner for industrial multi-robot use-case}{\small\href{https://doi.org/10.1109/ETFA54631.2023.10275420}{10.1109/ETFA54631.2023.10275420}}

\divider

\cvachievement{\faFile*[regular]}{Survey of communication-aware robot planning algorithms for 5G}{\small\href{https://doi.org/10.1109/ETFA52439.2022.9921449}{10.1109/ETFA52439.2022.9921449}}

\divider

\cvachievement{\faFile*[regular]}{Robotic measurement campaign \& Gaussian process model validation}{\small\href{https://doi.org/10.48550/arXiv.2603.08865}{10.48550/arXiv.2603.08865}}

\divider

\cvachievement{\faFile*[regular]}{Trustworthy edge computing for cyber-physical systems: a review}{\small\href{https://doi.org/10.1145/3539662}{10.1145/3539662}}

\hfill

\cvsection{References}

Available upon request.

\end{paracol}

%% ---- JOB LISTING -------------------------------------------------------


% AeroVect is transforming ground handling with autonomy, redefining how airlines and ground service providers around the globe run day-to-day operations. We are a Series A company backed by top-tier venture capital investors in aviation and autonomous driving. Our customers include some of the world’s largest airlines and ground handling providers. For more information, visit www.aerovect.com.  
  
% We are looking for an experienced Senior Software Engineer who can design and build best-in-class behavior planning systems for autonomous driving in structured, low-speed environments.  
  
% In this role, you'll own the design and implementation of key modules in the behavior planner — the decision-making layer that determines what the vehicle should do in complex, dynamic airside scenarios. You'll work at the intersection of mission-level goals and motion-level execution, tackling problems in multi-agent interaction modeling, rule-based and learned decision-making, and robust handling of edge cases unique to airport ground operations.  
  
% This opportunity offers a deeply technical engineer the chance to shape a market-defining enterprise product that combines autonomous vehicle technology with a robotics-as-a-service (RaaS) business model. This role reports to our Planning Tech Lead and works closely with the autonomy engineering team.  
  
% You Will:  

% - Develop and implement advanced behavior planning algorithms for autonomous vehicles
% - Collaborate with cross-functional teams to ensure robust integration and functionality of planning systems
% - Design, write, and maintain efficient and scalable code in C++ and Python
% - Contribute to the architecture and continuous improvement of behavior planning software
% - Conduct extensive testing in simulated environments and real-world scenarios to validate and refine behavior planning algorithms
% - Analyze system performance and implement enhancements based on data and feedback
% - Maintain comprehensive documentation of code, algorithms, and system designs
% - Work closely with other engineering teams to ensure seamless coordination and development

% You Have:  

% - Proficient in modern C++ (11/14/17) and object-oriented programming
% - Skilled in Python for rapid prototyping and testing
% - Strong in debugging, profiling, and optimizing code
% - Deep understanding of behavior planning algorithms such as state machines, behavior trees, and probabilistic planning
% - Familiarity with path planning algorithms like A\*, RRT, or optimization-based methods
% - Master’s degree in Computer Science, Robotics, or a related field
% - Minimum of 3 years of industry experience in autonomous driving, robotics, or a related field

% We Prefer:  

% - Knowledge of state machines, behavior trees, and decision-making under uncertainty
% - Expertise in path planning algorithms such as A\*, D\*, and Rapidly-exploring Random Trees (RRT)
% - Knowledge of machine learning techniques, especially in the context of behavior prediction and planning
% - Experience with ROS / ROS2
% - Implementing systems that can re-plan at high frequencies to adapt to dynamic changes in the environment
% - Ensuring that behavior planning algorithms can execute with minimal latency for real-time navigation
% - Proficiency in optimization techniques and probabilistic models for making informed planning decisions under uncertainty
% - Master’s degree or PhD in Robotics, AI, Mathematics, or a related field with a focus on planning, optimization, or control theory is a plus

%% ---- COVER LETTER -------------------------------------------------------

% AeroVect's behavior planner sits precisely at the intersection where my doctoral research has been conducted for the past five years: the boundary between mission-level goals and motion-level execution for fleets of autonomous vehicles operating in structured, dynamic environments. My thesis develops the theory and implementation of planners that must decide *what* a fleet of industrial robots should do and *when* — not just how they navigate — and that must remain robust as the environment and available resources change at runtime.

% A recurring theme in my research is that most planning work in this space reaches for a single monolithic MILP formulation and tries to fold everything — task sequencing, resource constraints, and trajectory feasibility — into one model. That approach gains tractability guarantees but sacrifices expressiveness: as the problem grows, the model becomes either intractable or so heavily approximated that the resulting plans are brittle. My work takes a different stance: use PDDL at the *mission level* to capture the full logical structure of goals, preconditions, and agent interactions, and keep motion-level execution as a separate, subordinate layer. This is a natural fit for a behavior planner, which is itself an intermediate layer in exactly such a hierarchy — responsible for what the vehicle should do in a given situation, with a motion planner below it and a mission or task manager above.

% I am not dogmatically attached to PDDL; I have worked with MILP, formal methods, and constraint-based formulations, as well as convex and nonlinear optimization, and my survey work has given me a systematic view of the strengths and failure modes of the main planning families. What I bring is the broader judgment about where to draw the abstraction boundaries — which representation to use at which level, and why combining them in a structured hierarchy outperforms trying to flatten everything into one formalism. That is precisely the kind of architectural thinking required to build a behavior planner that remains maintainable and extensible as edge cases accumulate.

% The airport apron is a different domain from an indoor factory, but the structural problem is a good match: a fleet of autonomous vehicles with interdependent goals must share a constrained physical space, respond to real-time changes, and maintain safe, predictable behavior throughout. I find this domain genuinely compelling — the combination of safety-critical operations, heterogeneous agent interactions, and tight scheduling constraints is a harder and more interesting version of the industrial logistics problem I have been studying.

% On the implementation side, I have years of C++ and Python development spanning bare-metal embedded systems, RTOS, and higher-level application software, as well as experience with ROS-based middleware from hands-on research such as the measurement campaign. I am comfortable working at the level of algorithm design, software architecture, and hardware-in-the-loop validation, and I have taught these topics at university level — which means I can communicate technical results clearly to colleagues from adjacent disciplines.

% I would welcome the opportunity to discuss how my background maps onto the specific challenges AeroVect is tackling.